\documentclass{article}
\usepackage[utf8]{inputenc}
\usepackage[T1]{fontenc}
\usepackage{geometry}
\usepackage{graphicx}
\usepackage{float}
\usepackage{hyperref}
\usepackage[english,indonesian]{babel}

\title{Optimasi Pemrosesan Template LaTeX Menggunakan Pendekatan Berbasis AI}
\author{Academic Support Agent}
\date{\today}

\begin{document}
\maketitle

\begin{abstract}

\end{abstract}

\chapter{PENDAHULUAN}


\chapter{TINJAUAN PUSTAKA}


% Original content (commented out):
% % Esta es un template de reporte de práctica en español para la materia de robótica de la coordinación de electrónica en la Universidad Iberoamericana
% Diseñado por Alexandro López González en enero de 2023

% Paquetes necesarios para el funcionamiento del documento-----------------------------------------------
\documentclass{article}
\usepackage[utf8]{inputenc}
\usepackage[spanish]{babel}
\decimalpoint
\usepackage{graphicx}
\graphicspath{ {Imagenes/} }
\usepackage[a4paper,width=150mm,top=25mm,bottom=25mm,bindingoffset=6mm]{geometry}
\usepackage{fancyhdr}
\pagestyle{fancy}
\usepackage{float}
\usepackage{amsmath}
\usepackage{caption}
\usepackage{subcaption}
\usepackage[bb=boondox,bbscaled=.95,cal=boondoxo]{mathalfa}
\usepackage{url}
\usepackage[hidelinks]{hyperref}
\usepackage{hyperref}
\usepackage{lscape}
% Fin de los paquetes ------------------------------------------------------------------------------------

\title{Reporte de práctica}
\author{Nombre apellidos \\ Nombre apellidos }

\begin{document} % Inicio del documento
% Título ------------------------------
\begin{figure}
    \centering
    \includegraphics{IBERO-IE.png}
\end{figure}
\maketitle
% fin del título ----------------------

\section{Resumen}

En esta sección se realiza una presentación concisa de la práctica, para describir de una manera sintetizada y/o abreviada el trabajo realizado. 

Sirve para sintetizar de manera objetiva y precisa el contenido y la organización de un documento el lector pueda tener una idea general de lo que se basa el trabajo.

En general se debe escribir al final del trabajo.

\subsection{Objetivos}
El objetivo general es... Los objetivos específicos son:
\begin{itemize}
    \item Objetivo 1
    \item Objetivo 2
    \item 
\end{itemize}

\section{Marco teórico}

Se detallan los conceptos teóricos y las ecuaciones, modelos matemáticos y referencias principales \cite{Spong2005-hq}.


\section{Procedimiento}
A modo de enumeración se pueden incluir los eventos realizados en el orden adecuado, detallando los procesos, herramientas, software y programas usados.

\begin{enumerate}
    \item 
    \item 
    \item 
\end{enumerate}


\section{Resultados}
Se presentan los resultados de las acciones tamadas, consecuencias medibles de los experimentos, planos de diseños u órdenes de compra.


\section{Conclusión}
Las conclusiones son una recopilación y análisis del proceso y resultado de la práctica, una conclusión no es una opinión o sugerencia, tampoco un resumen, sino un argumento lógicamente derivado de las premisas evaluadas, es decir, fruto de una labor de pensamiento y evaluación.



% Documentos de la bibliografía
\bibliography{Biblio.bib}
\bibliographystyle{plain}

\end{document}


\end{document}